% 2007 年四川卷理科第 2 题：f(x)=1+\log_2 x 与 g(x)=2^{1-x} 在同一坐标系下的四个选项图。
% 来源 PDF：gaokaomath-src/普通高考/2007/2007四川理.pdf，SHA256 6f3287a6ec1c136fb7e457cb0264d069c44eb8ff58a81cc3598f9c8f1395cc2d
% 原页：2007四川理.pdf 第 1 页（A3 三栏）第 1 栏，书页码 409，图在第 2 题题干下方。
% 重绘依据：四图为等大的独立直角坐标系（原点 O，两轴带箭头，两轴各标刻度 1、2），
%   (A)(B) 为上排、(C)(D) 为下排；(A) 递减曲线过 (0,2) 后向右渐近 x 轴、递增曲线在
%   x=1 处穿 x 轴，二者交于 x≈1.45；(B) 递减曲线过 (0,1)、递增曲线在 x=0.5 处穿
%   x 轴，二者交于 x≈0.75；(C) 为正确图象，直接按真函数采样 f(x)=1+\log_2 x 与
%   g(x)=2^{1-x}，二者交于 (1,1)；(D) 二者均递增，分别过 (-0.55,0) 与 (0.42,0)，
%   交于 x≈1.3。曲线只取原图画出的可见区间，不外推出坐标轴外。
\begin{tikzpicture}[scale=1.08, semithick, >=Stealth,
  every node/.style={inner sep=1.5pt}]
  % 四个等大的坐标系框（仅平移，画法完全一致）
  \foreach \dx/\dy/\opt in {0/0/A, 4.9/0/B, 0/-4.4/C, 4.9/-4.4/D}{
    \draw[->] ({\dx-0.65},\dy) -- ({\dx+2.85},\dy);
    \draw[->] (\dx,{\dy-0.7}) -- (\dx,{\dy+2.95});
    \node[anchor=north west] at ({\dx+1.07},{\dy-0.07}) {$1$};
    \node[anchor=north west] at ({\dx+2.07},{\dy-0.07}) {$2$};
    \node[anchor=east] at ({\dx-0.28},{\dy+1}) {$1$};
    \node[anchor=east] at ({\dx-0.28},{\dy+2}) {$2$};
    \node[anchor=north east] at ({\dx-0.08},{\dy-0.08}) {$O$};
    \node[anchor=west] at ({\dx+2.92},\dy) {$x$};
    \node[anchor=south west] at (\dx,{\dy+2.97}) {$y$};
    \node[anchor=east] at ({\dx-0.72},{\dy-0.56}) {(\opt)};
  }
  % (A) 递减过 (0,2)、递增过 (1,0)，交于 1<x<2
  \begin{scope}[shift={(0,0)}]
    \draw[domain=-0.35:2.45, samples=100, smooth]
      plot (\x,{2*exp(-0.65*\x)});
    \draw[domain=0.80:2.45, samples=100, smooth]
      plot (\x,{1.55*ln(\x)/ln(2)});
  \end{scope}
  % (B) 递减过 (0,1)、递增过 (0.5,0)，交于 x<1
  \begin{scope}[shift={(4.9,0)}]
    \draw[domain=-0.55:2.60, samples=100, smooth]
      plot (\x,{exp(-0.72*\x)});
    \draw[domain=0.42:2.15, samples=100, smooth]
      plot (\x,{ln(\x/0.5)/ln(2)});
  \end{scope}
  % (C) 真图象：f(x)=1+log_2 x 与 g(x)=2^{1-x}，交于 (1,1)
  \begin{scope}[shift={(0,-4.4)}]
    \draw[domain=-0.30:2.45, samples=100, smooth]
      plot (\x,{exp((1-\x)*ln(2))});
    \draw[domain=0.36:2.45, samples=100, smooth]
      plot (\x,{1+ln(\x)/ln(2)});
  \end{scope}
  % (D) 两条均递增，分别过 (-0.55,0)、(0.42,0)，交于 x≈1.3
  \begin{scope}[shift={(4.9,-4.4)}]
    \draw[domain=-0.64:2.50, samples=120, smooth]
      plot (\x,{1.35*ln(1+(\x+0.55)/1.1)});
    \draw[domain=0.18:1.90, samples=120, smooth]
      plot (\x,{1.60*sign(\x-0.42)*abs(\x-0.42)^1.25});
  \end{scope}
\end{tikzpicture}
